% =============================================================
%  ch14b_projekt0b_stereo.tex  --  Projekt 0.B: Stereo + Kalman
% =============================================================

\chapter{Projekt 0.B -- Stereokamera, Objekterkennung und Kalman-Tracking}
\label{ch:projekt0b}

\textbf{Ziel:} Eine simulierte \textbf{Stereokamera} beobachtet einen
\textbf{farbigen Ball}. Wir erkennen ihn im Bild (OpenCV), schätzen seine
\textbf{3D-Position und Entfernung} durch eigene \textbf{Triangulation} und
glätten bzw.\ prädizieren die Bahn mit einem selbst implementierten
\textbf{Kalman-Filter}. Das ist die direkte Vorübung zu Projekt~2
(Tischtennisroboter, Kapitel~\ref{ch:projekt0a} liefert den Bewegungsteil).

\begin{infobox}[title=Was wir übernehmen -- und was wir selbst bauen]
\textbf{Übernommen (Geometrie/Sensor):} ein Kugel-Primitive als Ball und der
Gazebo-Kamerasensor -- kein CAD nötig. \textbf{Selbst implementiert:} die
Farberkennung, die Stereo-Triangulation (Entfernung) und der Kalman-Filter.
Theorie dazu in den Kapiteln~\ref{ch:stereosehen} (Stereosehen) und
\ref{ch:kalman} (Kalman/EKF) -- hier setzen wir sie praktisch um.
\end{infobox}

\section{ROS\,2-Bausteine}
\begin{longtable}{@{}p{2.6cm}p{12cm}@{}}
\toprule
\textbf{Typ} & \textbf{Konkret (selbst implementiert, sofern nicht anders vermerkt)} \\
\midrule
Nodes    & Gazebo-Stereokamera + \code{ros\_gz}-Bildbrücke (Standard) $\cdot$
           \code{ball\_detector} (HSV-Erkennung L+R \emph{und} Triangulation) $\cdot$
           \code{ball\_tracker} (Kalman + Vorhersage) $\cdot$
           \code{marker\_viz} (RViz-Marker) \\
Topics   & \code{/stereo/left/image}, \code{/stereo/right/image} $\cdot$
           \code{/ball/point} (gemessen, \code{PointStamped}) $\cdot$
           \code{/ball/estimate} (KF) $\cdot$
           \code{/ball/prediction} (\code{nav\_msgs/Path}) $\cdot$
           \code{/ball/markers} (\code{MarkerArray}) $\cdot$ \code{/clock}, \code{/tf} \\
Services & \code{/tracker/reset} (KF zurücksetzen, \code{std\_srvs/Trigger}) \\
Parameter & HSV-Bereich (Ballfarbe) $\cdot$ Brennweite $f$, Basisbreite $B$,
            Hauptpunkt $(c_x,c_y)$ $\cdot$ KF-Rauschen $Q,R$, Vorhersage-Horizont \\
\bottomrule
\end{longtable}

\section{rqt\_graph (Zielbild)}
\begin{center}
\begin{tikzpicture}[node distance=5mm and 11mm]
\node[node] (cam) {/stereo\_camera\\\scriptsize (Gazebo+Bridge)};
\node[node,right=of cam] (det) {/ball\_detector\\\scriptsize (HSV L/R + Triangulation)};
\node[node,right=of det] (trk) {/ball\_tracker\\\scriptsize (Kalman + Vorhersage)};
\node[node,below=9mm of trk] (viz) {/marker\_viz};
\node[node,left=of viz] (rviz) {/rviz2};
\draw[flow] (cam) -- node[topic,fill=white,inner sep=1pt]{\scriptsize L/R image} (det);
\draw[flow] (det) -- node[topic,fill=white,inner sep=1pt]{\scriptsize ball/point} (trk);
\draw[flow] (trk) -- node[topic,fill=white,inner sep=1pt]{\scriptsize estimate/pred} (viz);
\draw[flow] (viz) -- node[topic,fill=white,inner sep=1pt]{\scriptsize markers} (rviz);
\draw[flow] (det) to[bend right=20] node[topic,fill=white,inner sep=1pt]{\scriptsize measured} (viz);
\end{tikzpicture}
\captionof{figure}{Projekt~0.B: Wahrnehmungskette von der Stereokamera über
Erkennung und Triangulation bis zum Kalman-Tracker und der RViz-Darstellung.}
\end{center}

\begin{figure}[htbp]\centering
  \includegraphics[width=0.95\textwidth]{{project_0.B/rqt_graph}.png}
  \caption{Tatsächlicher \code{rqt\_graph} von Projekt~0.B: Die
  \code{ros\_gz}-Bildbrücke speist \code{/stereo/left|right/image} in
  \code{ball\_detector}, dieser publiziert \code{/ball/point}; \code{ball\_tracker}
  erzeugt \code{/ball/estimate} und \code{/ball/prediction}, die \code{marker\_viz}
  und (über \code{fov\_viz}) RViz darstellen.}
  \label{fig:p0b-rqtgraph}
\end{figure}

\section{Schritt für Schritt}

\subsection{Schritt 0 -- Paket anlegen (ein Workspace für alle Projekte)}
Ein colcon-Workspace fasst beliebig viele Pakete; Projekt~0.B kommt deshalb als
weiteres Paket in dasselbe \code{\textasciitilde/ros2\_ws} neben die 0.A-Pakete --
\emph{kein} neuer Workspace nötig. \textbf{Lehre aus 0.A:} hier brauchen wir
\emph{keine} eigenen Interfaces; Standardnachrichten reichen
(\code{geometry\_msgs}, \code{nav\_msgs/Path}, \code{visualization\_msgs},
\code{std\_srvs}). Das erspart das gesamte
\code{rosidl}/\code{ament\_cmake}-Drumherum.
\begin{lstlisting}[style=bash,caption={Paket anlegen}]
cd ~/ros2_ws/src
ros2 pkg create --build-type ament_python p0b_stereo --license Apache-2.0 \
  --dependencies rclpy sensor_msgs geometry_msgs nav_msgs \
                 visualization_msgs std_srvs cv_bridge message_filters
mkdir -p p0b_stereo/{worlds,launch,config,p0b_stereo}
\end{lstlisting}
\begin{lstlisting}[style=bash,caption={\code{ros2\_ws}-Delta nach Schritt 0 (\# neu)}]
src/p0b_stereo/   # neu -- ament_python, Standard-Messages (keine Interfaces!)
\end{lstlisting}

\subsection{Schritt 1 -- Welt: Ball und Stereokamera}
Zwei Kamerasensoren im Abstand der Basisbreite $B=0{,}10\,$m, Blick entlang
$+x$, und ein orangefarbener Ball (dynamisch -- er fällt und lässt sich in der
GUI mit der Maus bewegen). Wichtig: \code{<optical\_frame\_id>} setzt den Rahmen,
in dem wir später die 3D-Punkte veröffentlichen.
\begin{lstlisting}[style=xml,caption={\code{worlds/stereo.sdf} (Auszug -- Ground/Sun/Plugins analog)}]
<model name="stereo_rig"><static>true</static><pose>0 0 1 0 0 0</pose>
  <link name="base">
    <visual name="body">          <!-- sichtbarer Koerper des Rigs in Gazebo -->
      <geometry><box><size>0.05 0.18 0.05</size></box></geometry></visual>
    <sensor name="left" type="camera">
      <pose>0 0.05 0 0 0 0</pose>           <!-- +B/2 -->
      <topic>/stereo/left/image</topic>
      <update_rate>30</update_rate><always_on>1</always_on>
      <camera><horizontal_fov>1.047</horizontal_fov>
        <image><width>640</width><height>480</height></image>
        <optical_frame_id>stereo_left_optical</optical_frame_id></camera>
    </sensor>
    <sensor name="right" type="camera">
      <pose>0 -0.05 0 0 0 0</pose>          <!-- -B/2 -->
      <topic>/stereo/right/image</topic>
      <update_rate>30</update_rate><always_on>1</always_on>
      <camera><horizontal_fov>1.047</horizontal_fov>
        <image><width>640</width><height>480</height></image>
        <optical_frame_id>stereo_right_optical</optical_frame_id></camera>
    </sensor>
  </link>
</model>

<model name="ball"><pose>3 0 1.5 0 0 0</pose>
  <link name="link">
    <inertial><mass>0.05</mass><inertia><ixx>7.2e-5</ixx><iyy>7.2e-5</iyy>
      <izz>7.2e-5</izz><ixy>0</ixy><ixz>0</ixz><iyz>0</iyz></inertia></inertial>
    <collision name="c"><geometry><sphere><radius>0.06</radius></sphere></geometry></collision>
    <visual name="v"><geometry><sphere><radius>0.06</radius></sphere></geometry>
      <material><ambient>1 0.25 0 1</ambient><diffuse>1 0.25 0 1</diffuse></material></visual>
  </link>
</model>
\end{lstlisting}
\noindent Die vollständige Datei (mit \code{ground}, \code{sun} und den
Gazebo-\code{system}-Plugins inkl.\ \code{Sensors}/\code{ogre2}) liegt in
\code{worlds/stereo.sdf}; gestartet wird sie über die Launch-Datei aus
Schritt~6.
\begin{lstlisting}[style=bash,caption={\code{ros2\_ws}-Delta nach Schritt 1 (\# neu)}]
src/p0b_stereo/worlds/stereo.sdf   # neu (Stereo-Rig + Ball + Welt)
\end{lstlisting}

\begin{infobox}[title={SDF (0.B) vs.\ URDF/Xacro (0.A) -- der Unterschied}]
\textbf{URDF/Xacro} (Projekt~0.A, \code{arm.urdf.xacro}) beschreibt \emph{einen
Roboter} als kinematischen Baum aus Links und Joints. Es ist ein \emph{ROS}-Format:
\code{robot\_state\_publisher}, \code{ros2\_control} und RViz lesen es direkt;
Xacro ist nur ein Makro-Präprozessor darüber. Sensoren und Physik müssen über
\code{<gazebo>}-Erweiterungen angehängt werden.

\textbf{SDF} (Projekt~0.B, \code{stereo.sdf}) ist das \emph{native Format von
Gazebo} und beschreibt eine ganze \emph{Welt}: Licht, Boden, mehrere Modelle,
\textbf{Sensoren} (Kamera, IMU, LiDAR) und \textbf{System-Plugins} (Physik,
Rendering) -- alles erstklassig. ROS kennt SDF nicht direkt; die Verbindung
entsteht über \code{ros\_gz}-Brücken.

\textbf{Faustregel:} \emph{einen Roboter} für ROS beschreiben $\rightarrow$
URDF/Xacro (wird beim Spawnen intern nach SDF gewandelt). \emph{Eine
Simulationswelt mit Sensoren} bauen $\rightarrow$ SDF. In 0.A war der Arm der
Gegenstand (URDF, in Gazebo gespawnt); in 0.B ist die Kamera-Welt der Gegenstand
(SDF), und die Wahrnehmungsknoten leben rein in ROS.
\end{infobox}

\subsection{Schritt 2 -- Bilder und Zeit nach ROS bruecken}
Die Bilder bringt die \code{ros\_gz}-Bildbrücke nach ROS; die Simulationszeit
(\code{/clock}) braucht eine eigene Brücke -- \textbf{exakt die Lehre aus 0.A}:
ohne \code{/clock} stehen alle \code{use\_sim\_time}-Knoten ohne Uhr da.
\begin{lstlisting}[style=bash,caption={Brücken (in der Launch-Datei aus Schritt 6 gebündelt)}]
ros2 run ros_gz_image image_bridge /stereo/left/image /stereo/right/image
ros2 run ros_gz_bridge parameter_bridge /clock@rosgraph_msgs/msg/Clock[gz.msgs.Clock
\end{lstlisting}

\subsection{Schritt 3 -- Objekterkennung (OpenCV, HSV)}
In beiden Bildern den Ball über seine Farbe finden und den Pixelmittelpunkt
$(u,v)$ bestimmen (Kern des \code{ball\_detector}):
\begin{lstlisting}[style=py,caption={HSV-Erkennung (Methode \code{detect})}]
def detect(self, img_msg):
    img = self.bridge.imgmsg_to_cv2(img_msg, 'bgr8')
    hsv = cv2.cvtColor(img, cv2.COLOR_BGR2HSV)
    mask = cv2.inRange(hsv, self.lo, self.hi)     # self.lo/hi = HSV-Grenzen
    cnts, _ = cv2.findContours(mask, cv2.RETR_EXTERNAL, cv2.CHAIN_APPROX_SIMPLE)
    if not cnts:
        return None
    (u, v), r = cv2.minEnclosingCircle(max(cnts, key=cv2.contourArea))
    return (float(u), float(v)) if r >= 3.0 else None
\end{lstlisting}

\subsection{Schritt 4 -- Triangulation und der vollständige Detektor}
Bei achsparalleler Stereoanordnung ist die Disparität $d=|u_L-u_R|$; mit
Brennweite $f$ (Pixel) und Basisbreite $B$:
\begin{align*}
Z=\frac{f\,B}{d},\qquad X=\frac{(u_L-c_x)\,Z}{f},\qquad Y=\frac{(v_L-c_y)\,Z}{f}.
\end{align*}
$f,c_x,c_y$ berechnen wir aus den bekannten Kameraparametern
($f=\tfrac{w/2}{\tan(\text{hfov}/2)}$) -- so ist \emph{keine}
\code{camera\_info}-Brücke nötig (eine Fehlerquelle weniger). Beide Bilder werden
per \code{ApproximateTimeSynchronizer} zum selben Zeitpunkt zusammengeführt:
\begin{lstlisting}[style=py,caption={\code{p0b\_stereo/ball\_detector.py} (Kern)}]
class BallDetector(Node):
    def __init__(self):
        super().__init__('ball_detector')
        w, h, hfov = 640, 480, 1.047
        self.f = (w/2.0) / np.tan(hfov/2.0)          # Brennweite [px]
        self.cx, self.cy, self.B = w/2.0, h/2.0, 0.10
        self.frame_id = 'stereo_left_optical'
        self.lo = np.array([3, 120, 70]); self.hi = np.array([25, 255, 255])
        self.bridge = CvBridge()
        self.pub = self.create_publisher(PointStamped, '/ball/point', 10)
        left  = Subscriber(self, Image, '/stereo/left/image')
        right = Subscriber(self, Image, '/stereo/right/image')
        self.sync = ApproximateTimeSynchronizer([left, right], 10, 0.05)
        self.sync.registerCallback(self.on_pair)

    def on_pair(self, left_msg, right_msg):
        L, R = self.detect(left_msg), self.detect(right_msg)
        if L is None or R is None: return
        (uL, vL), (uR, _) = L, R
        d = abs(uL - uR)
        if d < 1e-3: return                          # Objekt zu weit / kein Match
        Z = self.f * self.B / d
        X = (uL - self.cx) * Z / self.f
        Y = (vL - self.cy) * Z / self.f
        msg = PointStamped()
        msg.header.stamp = left_msg.header.stamp
        msg.header.frame_id = self.frame_id
        msg.point.x, msg.point.y, msg.point.z = float(X), float(Y), float(Z)
        self.pub.publish(msg)
\end{lstlisting}
\noindent Im echten Knoten sind HSV, $B$ und Brennweite über
\code{declare\_parameter} einstellbar. Vertiefung: Kapitel~\ref{ch:stereosehen},
\ref{ch:kalibrierung}.
\begin{lstlisting}[style=bash,caption={\code{ros2\_ws}-Delta nach Schritt 4 (\# neu / \# angepasst)}]
src/p0b_stereo/
|- setup.py                     # angepasst: entry_points -> ball_detector
`- p0b_stereo/ball_detector.py  # neu (Erkennung + Triangulation)
\end{lstlisting}

\subsection{Schritt 5 -- Kalman-Tracker mit Vorhersage}
Lineares Modell konstanter Geschwindigkeit, Zustand
$\mathbf{x}=[X,Y,Z,\dot X,\dot Y,\dot Z]^\top$:
\begin{align*}
\hat{\mathbf{x}}^- = F\hat{\mathbf{x}},\ P^-=FPF^\top+Q;\qquad
K=P^-H^\top(HP^-H^\top+R)^{-1},\ \hat{\mathbf{x}}=\hat{\mathbf{x}}^-+K(\mathbf z-H\hat{\mathbf x}^-).
\end{align*}
\begin{lstlisting}[style=py,caption={\code{p0b\_stereo/ball\_tracker.py} (Kern)}]
def on_point(self, msg):
    t = msg.header.stamp.sec + msg.header.stamp.nanosec*1e-9
    z = np.array([msg.point.x, msg.point.y, msg.point.z])
    if not self.initialized:
        self.x[:3] = z; self.initialized = True; self.last_t = t; return
    dt = max(1e-3, t - self.last_t); self.last_t = t
    F = self._F(dt)
    self.x = F @ self.x                                # Praediktion
    self.P = F @ self.P @ F.T + self.Q
    S = self.H @ self.P @ self.H.T + self.R            # Korrektur
    K = self.P @ self.H.T @ np.linalg.inv(S)
    self.x = self.x + K @ (z - self.H @ self.x)
    self.P = (np.eye(6) - K @ self.H) @ self.P
    self._publish(msg.header.stamp, dt)
\end{lstlisting}
\noindent \code{\_publish} sendet \code{/ball/estimate} und integriert den
Zustand \emph{ohne} Messung $N$ Schritte vorwärts -- das ergibt
\code{/ball/prediction} (\code{nav\_msgs/Path}), die Grundlage fürs Abfangen in
Projekt~2. Ein \code{/tracker/reset} (\code{std\_srvs/Trigger}) nullt den Filter.
\begin{lstlisting}[style=bash,caption={\code{ros2\_ws}-Delta nach Schritt 5 (\# neu / \# angepasst)}]
src/p0b_stereo/
|- setup.py                    # angepasst: entry_points -> ball_tracker
`- p0b_stereo/ball_tracker.py  # neu (Kalman + Vorhersage + reset-Service)
\end{lstlisting}

\subsection{Schritt 6 -- Marker in RViz und alles starten}
\code{marker\_viz} macht Messung (rot), Schätzung (grün) und Vorhersage (blaue
Linie) als \code{MarkerArray} sichtbar:
\begin{lstlisting}[style=py,caption={\code{p0b\_stereo/marker\_viz.py} (Kern)}]
def on_meas(self, msg):   # rote Kugel an der gemessenen Position
    self._emit(self._sphere('measured', msg.point, (1, 0.1, 0),
                            msg.header.frame_id, msg.header.stamp, 0.08))

def on_pred(self, msg):   # blaue Linie entlang der Vorhersage
    m = Marker(); m.header = msg.header; m.ns = 'prediction'
    m.type = Marker.LINE_STRIP; m.action = Marker.ADD
    m.scale.x = 0.02; m.color.b = 1.0; m.color.a = 1.0
    m.points = [ps.pose.position for ps in msg.poses]
    self._emit(m)
\end{lstlisting}
\noindent Eine Launch-Datei startet alles in \emph{einem} Befehl -- Gazebo,
\code{/clock}-Brücke, Bildbrücke, eine statische TF (damit der Kamerarahmen in
RViz existiert -- Lehre aus 0.A), die drei Knoten und RViz, alle mit
\code{use\_sim\_time}:
\begin{lstlisting}[style=py,caption={\code{launch/stereo.launch.py} (Knotenliste)}]
return LaunchDescription([
    gz,             # ros_gz_sim + worlds/stereo.sdf  (gz_args: '-r')
    clock_bridge,   # /clock  (gz -> ROS)
    image_bridge,   # /stereo/left|right/image
    static_tf,      # world -> stereo_left_optical
    detector, tracker, viz,   # parameters=[{'use_sim_time': True}]
    rviz,           # -d config/stereo.rviz
])
\end{lstlisting}
\begin{lstlisting}[style=bash,caption={Terminal 1 -- bauen, sourcen, alles starten}]
cd ~/ros2_ws
colcon build --packages-select p0b_stereo
source install/setup.bash
ros2 launch p0b_stereo stereo.launch.py
\end{lstlisting}
\noindent Die Launch-Datei startet \emph{auch} die beiden Brücken aus Schritt~2
(Bild + \code{/clock}). Startet man die Knoten \emph{einzeln} (z.\,B.\ zum
Debuggen), müssen diese Brücken in einem eigenen Terminal laufen -- aber
\textbf{nicht zusätzlich} zur Launch, sonst doppelte Publisher:
\begin{lstlisting}[style=bash,caption={Brücken (nur bei manuellem Start nötig)}]
ros2 run ros_gz_image image_bridge /stereo/left/image /stereo/right/image
ros2 run ros_gz_bridge parameter_bridge /clock@rosgraph_msgs/msg/Clock[gz.msgs.Clock
\end{lstlisting}
\begin{lstlisting}[style=bash,caption={Terminal 2 -- mitlesen / Filter zuruecksetzen}]
source ~/ros2_ws/install/setup.bash
ros2 topic echo /ball/point                        # gemessene 3D-Position
ros2 service call /tracker/reset std_srvs/srv/Trigger {}
\end{lstlisting}
\noindent \textbf{Tipp:} Den Ball in Gazebo mit der Maus greifen und bewegen --
Messung, Schätzung und Vorhersage folgen in RViz. \textbf{Sichtfeld im Blick
behalten:} RViz zeigt die beiden Kamerabilder (\code{/stereo/left|right/image});
der Ball wird nur verfolgt, solange er in \emph{beiden} sichtbar ist. Verlässt er
das Bild, erkennt \code{ball\_detector} nichts mehr und die Marker bleiben auf der
letzten Position stehen -- kein Fehler. Das schwarze Rig-Kästchen zeigt die
Kameraposition in Gazebo.
\begin{figure}[htbp]\centering
  \includegraphics[width=0.95\textwidth]{{project_0.B/step6_rviz_gazebo_terminal}.png}
  \caption{Projekt~0.B in Betrieb: links Gazebo mit Stereo-Rig und Ball, rechts
  RViz mit den Markern (gemessen rot, Kalman-Schätzung grün, Vorhersage blau) und
  den beiden Kamerabildern; im Terminal die \code{/ball/point}-Ausgabe. Der Filter
  glättet das Messrauschen und schätzt voraus.}
  \label{fig:p0b-rviz-markers}
\end{figure}
\begin{lstlisting}[style=bash,caption={\code{ros2\_ws}-Delta nach Schritt 6 (\# neu / \# angepasst)}]
src/p0b_stereo/
|- setup.py                  # angepasst: entry_points -> marker_viz, data_files
|- p0b_stereo/marker_viz.py  # neu (MarkerArray fuer RViz)
|- launch/stereo.launch.py   # neu (alles in einem Start)
`- config/stereo.rviz        # neu (Fixed Frame, MarkerArray, Image)
\end{lstlisting}

\section{Implementierungs-Referenz: Parameter und Synchronisation}
Die rclpy-Grundlagen (Publisher, Subscription, Service, Launch) stehen in der
Referenz von Projekt~0.A. In 0.B kommen zwei neue Bausteine dazu.

\subsection{Parameter -- \code{declare\_parameter} / \code{get\_parameter}}
Statt Werte (HSV-Bereich, Basisbreite, Brennweite) fest zu verdrahten, deklariert
man sie als \emph{Parameter}: so sind sie beim Start oder zur Laufzeit änderbar,
ohne den Code anzufassen.
\begin{lstlisting}[style=py]
self.declare_parameter('baseline', 0.10)        # Name + Default(-typ)
B = self.get_parameter('baseline').value        # Wert auslesen
\end{lstlisting}
\noindent Gesetzt wird z.\,B.\ in der Launch-Datei
(\code{parameters=[\{'baseline': 0.12\}]}) oder per CLI
(\code{--ros-args -p baseline:=0.12}). Tutorial:
\href{https://docs.ros.org/en/jazzy/Tutorials/Beginner-Client-Libraries/Using-Parameters-In-A-Class-Python.html}{docs.ros.org -- Parameters in a Class (Python)}.

\subsection{Zwei Topics zeitlich koppeln -- \code{message\_filters}}
Die Triangulation braucht das linke \emph{und} rechte Bild \emph{desselben}
Moments. Ein gewöhnliches \code{create\_subscription} pro Kamera liefert die
Bilder unabhängig -- man wüsste nicht, welche zusammengehören. Der
\code{ApproximateTimeSynchronizer} verknüpft beide Subscriptions und ruft den
Callback erst, wenn ein zeitlich passendes \emph{Paar} vorliegt:
\begin{lstlisting}[style=py]
from message_filters import Subscriber, ApproximateTimeSynchronizer
left  = Subscriber(self, Image, '/stereo/left/image')
right = Subscriber(self, Image, '/stereo/right/image')
sync = ApproximateTimeSynchronizer([left, right], queue_size=10, slop=0.05)
sync.registerCallback(self.on_pair)   # on_pair(left_msg, right_msg)
\end{lstlisting}
\noindent \code{slop} ist die maximal erlaubte Zeitdifferenz (s) zwischen den
beiden Zeitstempeln. Doku:
\href{https://docs.ros.org/en/jazzy/p/message_filters/}{docs.ros.org -- message\_filters}.

\section{Sichtfeld visualisieren und Sim-to-Real}

\subsection{Wie das Sichtfeld definiert ist}
Jede Kamera in der SDF hat \code{<horizontal\_fov>} (horizontaler Öffnungswinkel)
und \code{<image><width/height>}. Daraus folgt die gesamte Abbildungsgeometrie:
\begin{itemize}
\item \textbf{Brennweite (Pixel):} $f=\dfrac{w/2}{\tan(\text{hfov}/2)}$ -- genau die
  Größe, die \code{ball\_detector} zur Triangulation berechnet
  ($640\,$px, $60^\circ \Rightarrow f\approx 554$).
\item \textbf{Vertikaler Öffnungswinkel} über das Seitenverhältnis:
  $\text{vfov}=2\arctan\!\big(\tan(\text{hfov}/2)\cdot h/w\big)$.
\item \textbf{Hauptpunkt} $(c_x,c_y)=(w/2,\,h/2)$ -- ideale Kamera, ohne Verzeichnung.
\end{itemize}
Bisher \emph{sehen} wir das Sichtfeld nur indirekt über die beiden Kamerabilder.

\subsection{Das Frustum als Pyramide zeigen (\code{fov\_viz})}
Der Knoten \code{fov\_viz} macht das Sichtfeld im 3D sichtbar: pro Kamera eine
Pyramide (Apex am Kamerazentrum, Grundfläche auf der Fernebene) als
\code{LINE\_LIST}-Marker auf \code{/stereo/fov\_markers}. Die Ecken folgen direkt
aus hfov, Seitenverhältnis und Reichweite:
\begin{lstlisting}[style=py,caption={\code{fov\_viz.py} (Kern der Frustum-Berechnung)}]
tan_h = np.tan(hfov/2.0); tan_v = tan_h * h/w          # Halbwinkel
X, Y = far*tan_h, far*tan_v                            # halbe Fernflaeche
corners = [Point(x=apex_x + sx*X, y=sy*Y, z=far)       # optisch: x rechts, y unten
           for sx, sy in [(-1,-1), (1,-1), (1,1), (-1,1)]]
# LINE_LIST: Apex->4 Ecken + Fernrechteck; linke Kamera apex_x=0, rechte apex_x=B
\end{lstlisting}
\noindent In RViz erscheinen zwei überlappende Pyramiden (grün = links, blau =
rechts). Ihr \emph{Schnittvolumen} ist der Bereich, in dem \emph{beide} Kameras
den Ball sehen -- nur dort kann trianguliert werden. \code{fov\_viz} ist in
\code{stereo.launch.py} enthalten und als \emph{CameraFOV}-Display in
\code{stereo.rviz} eingebunden.
\begin{lstlisting}[style=bash,caption={\code{ros2\_ws}-Delta (\# neu / \# angepasst)}]
src/p0b_stereo/
|- p0b_stereo/fov_viz.py     # neu (Frustum-Marker beider Kameras)
|- setup.py                  # angepasst: entry_points -> fov_viz
|- launch/stereo.launch.py   # angepasst: fov_viz gestartet
`- config/stereo.rviz        # angepasst: CameraFOV- + RightImage-Display
\end{lstlisting}

\begin{figure}[htbp]\centering
  \includegraphics[width=0.95\textwidth]{{project_0.B/field_of_view}.png}
  \caption{Die von \code{fov\_viz} gezeichneten Sichtfeld-Pyramiden beider Kameras
  in RViz (grün = links, blau = rechts). Das Überlappungsvolumen ist der Bereich,
  in dem \emph{beide} Kameras den Ball sehen -- nur dort liefert die Triangulation
  eine 3D-Position.}
  \label{fig:p0b-fov}
\end{figure}
Damit das in Simulation Entwickelte später auf einer echten USB-Stereokamera
(z.\,B.\ einem ELP-Dual-Lens-Modul mit Global-Shutter) läuft, gleicht man drei
Größen mit dem realen Gerät ab -- dann teilen Sim und Realität dieselbe Geometrie:
\begin{enumerate}
\item \textbf{Öffnungswinkel/Brennweite:} \code{<horizontal\_fov>} auf den Wert des
  realen Objektivs setzen (Datenblatt oder Kalibrierung,
  Kapitel~\ref{ch:kalibrierung}) -- damit stimmt $f$.
\item \textbf{Auflösung:} \code{<image><width/height>} auf die reale
  Sensorauflösung pro Auge.
\item \textbf{Basisbreite $B$:} den Kameraabstand in der SDF (\code{<pose>}) und
  den \code{baseline}-Parameter auf den realen Linsenabstand des Moduls.
\end{enumerate}
\begin{infobox}[title=Was die Simulation (noch) nicht abbildet]
Reale Module haben \textbf{Linsenverzeichnung}: Gazebo kann sie über
\code{<distortion>} (Brown-Conrady $k_1,k_2,\ldots$) nachbilden -- sauberer ist,
\emph{beide} Seiten zu kalibrieren und die Bilder zu entzerren
(Kapitel~\ref{ch:kalibrierung}). Ein \textbf{Global-Shutter} (wie beim genannten
Modul) passt gut zum einfachen Pinhole-Modell; ein Rolling-Shutter erzeugt bei
schnellen Bällen Schräglagen, die wir hier nicht simulieren. Auch reale
Belichtung und Bildrauschen fehlen -- die HSV-Schwellen müssen am Ende an echten
Bildern nachjustiert werden.
\end{infobox}

\section{Häufige Fehler und Checkliste (Projekt 0.B)}
\begin{warnbox}[title={Stolpersteine -- vieles aus 0.A übertragbar}]
\begin{itemize}
\item \textbf{\code{/clock} nicht gebrückt / \code{use\_sim\_time} fehlt.} Wie in
  0.A: ohne \code{/clock}-Brücke und \code{use\_sim\_time:=true} fehlt die
  gemeinsame Zeit -- RViz zeigt nichts oder ``No transform''.
\item \textbf{Kamerarahmen existiert nicht in RViz.} \code{<optical\_frame\_id>}
  in der SDF setzen \emph{und} eine statische TF (\code{world\,$\rightarrow$\,%
  stereo\_left\_optical}) publizieren -- sonst ``Fixed Frame does not exist''.
\item \textbf{Kein Kamerabild (EGL/Headless).} Die Gazebo-\code{Sensors} brauchen
  einen funktionierenden Renderer (\code{ogre2}); auf der RTX\,3060 hilft die
  EGL-Einstellung aus dem Simulationskapitel -- sonst bleiben die Bild-Topics leer.
\item \textbf{HSV-Bereich zu eng/weit.} Kein Blob $\Rightarrow$ keine Punkte. Mit
  \code{rqt\_image\_view} und der Maske abstimmen; Beleuchtung der Welt beachten.
\item \textbf{Disparität $\le 0$.} Bei vertauschter Links/Rechts-Zuordnung wird
  $d$ negativ; wir nehmen $|u_L-u_R|$ und verwerfen $d\approx 0$ (Objekt zu weit).
\item \textbf{Bilder nicht synchron.} Ohne \code{ApproximateTimeSynchronizer}
  trianguliert man Links/Rechts aus verschiedenen Frames -- springende Tiefe.
\item \textbf{NumPy-Skalare in ROS-Felder.} \code{msg.point.x} erwartet ein
  Python-\code{float} -- mit \code{float(\ldots)} casten.
\item \textbf{Unnötige Custom-Interfaces.} Standardnachrichten nutzen
  (\code{geometry\_msgs}, \code{nav\_msgs/Path}, \code{visualization\_msgs}) --
  spart das \code{rosidl}-Drumherum aus 0.A.
\item \textbf{``Stereo is NOT SUPPORTED'' von RViz.} Harmlos -- gemeint ist
  OpenGL-3D-Brillen-Stereo, nicht die Stereokamera. Ignorieren.
\item \textbf{Marker stehen still / kein \code{/ball/point} mehr.} Der Ball ist
  nicht in \emph{beiden} Kamerabildern. Die Image-Panels in RViz beobachten und
  den Ball im gemeinsamen Sichtfeld halten -- erwartetes Verhalten, kein Bug.
\end{itemize}
\end{warnbox}
\begin{infobox}[title=Schnelle Prüfbefehle]
\begin{lstlisting}[style=bash]
ros2 topic hz /stereo/left/image     # liefert die Kamera Bilder?
ros2 topic echo /ball/point          # wird erkannt + trianguliert?
ros2 topic echo /ball/estimate       # schaetzt der Kalman-Filter?
rqt_graph                            # haengt die Kette zusammen?
\end{lstlisting}
\end{infobox}

\section{Was du danach beherrschst}
\begin{itemize}
\item \textbf{Sensorik in Simulation:} Gazebo-Kamerabilder über \code{ros\_gz}
  nach ROS\,2 bringen.
\item \textbf{Bildverarbeitung:} ein Objekt per OpenCV/\code{cv\_bridge} robust
  erkennen.
\item \textbf{Stereo-Geometrie:} Entfernung aus Disparität selbst triangulieren.
\item \textbf{Zustandsschätzung:} einen Kalman-Filter implementieren, der glättet
  \emph{und} prädiziert (Kapitel~\ref{ch:kalman}).
\item \textbf{Visualisierung:} Mess- und Schätzgrößen als RViz-Marker prüfen --
  die direkte Vorstufe zu Projekt~2.
\end{itemize}
